\section*{Предисловие}
\addcontentsline{toc}{section}{Предисловие}
\label{sec:predislovie}

Курсовият проект е разработен по дисциплината „Процесуално-правни аспекти на
киберпрестъпленията“ в Техническия университет~-- София. Темата „Ролята на прокурора в
досъдебната фаза на наказателния процес“ е възложена от научния ръководител
гл.~ас. д-р Боян Бележков.

Темата е тясно свързана с предмета на дисциплината. При киберпрестъпленията изходът от
делото обикновено се решава още в досъдебното производство.
Данните, които доказват престъплението, се съхраняват за кратко, лесно се променят и
често се намират у доставчици в други държави. Затова е важно кой и колко бързо може да
поиска запазването им, кой разрешава достъпа до тях и кой отговаря за законосъобразното
им събиране. В българския наказателен процес отговорът на тези въпроси минава през
фигурата на прокурора.

Работата е написана от студент с техническа подготовка. Затова в нея са разгледани не
само правомощията на прокурора по закон, но и начинът, по който те се упражняват по отношение на
компютърните данни, данните за трафика и криптоактивите. Нормативните актове са
цитирани по консолидираните им текстове към 16~ноември 2026~г. Последното отразено
изменение на Наказателно-процесуалния кодекс е обнародвано в ДВ, бр.~59 от 2026~г.

Авторът благодари на гл.~ас. д-р Боян Бележков за насоките при разработването на темата.
